On éclaire une fente verticale de largeur $a=\qty{0.1}{\mm}$, à l'aide d'un
laser qui donne une lumière monochromatique de longueur d'onde
$\lambda=\qty{632.8}{\nm}$. On observe sur un écran placé à la distance $D$ de
la fente, des taches lumineuses mettant en évidence le phénomène de diffraction.
La largeur de la tache centrale s'exprime par~: $\mathrm{L}=\frac{2\lambda
D}{a}$. La célérité de la lumière dans le vide (ou l'air) est
$c=\qty{3e8}{\m\per\s}$.

\begin{enumerate}
    \item Déterminer la valeur de la fréquence $v$ de la lumière utilisée.
    \item On refait l'expérience en utilisant un fil très fin vertical de
          diamètre $a_{0}$, on obtient une tache centrale de largeur $L_{0}=2L$.
          Déterminer la valeur de $a_{0}$.
\end{enumerate}

